\chapter{Epistemological and Methodological Framework}

\epigraph{\emph{Knowledge begins where translation becomes explicit.}}{}

The previous chapter assembled a distributed convergence. It showed that multiple literatures already encounter structured relations among frequencies, modes, recurrences, and resonant organizations, even when they name them differently and study them at different scales. That map, however, remains incomplete until one asks a harder question: what kind of knowledge could legitimately be produced about such an object? A convergent literature does not automatically yield a convergent method. It does not by itself tell us how to compare evidence across domains, how to organize claims of unequal strength, or how to decide what counts as support, limitation, or falsification when the object under study is partly physical, partly biological, partly experiential, and partly computational.

This chapter addresses that problem directly. Its task is to define the epistemological and methodological framework through which HIT proposes to transform convergence into evidence. The argument proceeds from broad positioning to operational structure. It first clarifies what kind of knowledge HIT seeks to produce and why that knowledge cannot be reduced either to laboratory science alone or to unconstrained experiential testimony. It then examines the institutional and representational limits that shape formalization, before turning to the specific roles played by computational modeling, the distinction between natural and perceptual harmony, Beacon (the program's experiential branch, developed more fully in Chapter 12), and the architecture of evidence itself. The chapter therefore functions as a hinge. It does not replace the empirical map with a manifesto about knowledge, nor does it collapse into a technical protocol manual. It establishes the conditions under which the claims of this manuscript can be made, tested, and revised with discipline.

\section{What kind of knowledge HIT produces}

HIT operates at the intersection of three traditions that are often kept apart: the natural sciences, practice-based research, and critical epistemology. From the sciences it inherits measurement, comparison, falsification, and the obligation to distinguish observation from inference. From practice-based inquiry it inherits the recognition that some phenomena first become legible through situated action, bodily adjustment, and disciplined attention rather than through prior formal theory alone \parencite{borgdorff_2012, smith_dean_2009}. From critical epistemology it inherits reflexivity about perspective, institution, and the conditions under which some forms of evidence become legible while others remain marginal \parencite{haraway_1988, harding_1991}. HIT is not a compromise among these traditions so much as an attempt to take seriously the fact that its object already exceeds the jurisdiction of any one of them taken alone.

This is why the question of truth must be handled with some care at the outset. Nietzsche's account of truth in \emph{The Gay Science} is useful here, not because it dissolves rigor, but because it reminds us that knowledge is never produced from nowhere and that human concepts have historically been shaped by forms of practical and vital orientation rather than by pure contemplative transparency \parencite{nietzsche_1882}. From the perspective adopted in this manuscript, scientific truth is not access to the Real. It is a highly disciplined way of constructing reality within the symbolic procedures of science: measurement, formalization, comparison, and controlled inscription. The Real, in the stronger sense that becomes clearer through Lacan and Derrida, is precisely what exceeds complete symbolic capture and returns only as resistance, bodily event, remainder, or rupture. HIT therefore does not oppose science to truth, but it refuses to treat scientific discourse as a transparent passage to what lies outside discourse. Different investigative regimes construct different realities, and a mature research program must learn how to compare and articulate those constructions without mistaking any one of them for final possession of the Real.

Maturana's reflexive formulation sharpens this point from within biology itself. If everything said is said by an observer, then knowledge can never be separated from the living system that produces the distinctions through which its world becomes articulate \parencite{maturana_1970, maturana_1988}. This is not an anti-scientific claim. It is a biologically grounded version of the same refusal of view-from-nowhere that animates Haraway's situated knowledges. Varela's neurophenomenology extends the consequence methodologically by arguing that first-person experience is not a contaminant to be purged from inquiry, but a domain of disciplined data that must be placed in rigorous relation with neurobiological description \parencite{varela_1996_neuro}. For HIT this convergence matters because it protects the framework from a false choice between formal science and lived evidence. Reflexivity is not a concession to subjectivism. It is part of the epistemic honesty demanded by any program that studies embodied, resonant, and partially self-interpreting systems.

The practical importance of this position becomes clear in the case of Beacon, the program's experiential and physiological probe, treated more fully in Chapter 12. If participants repeatedly report shifts in regulation, resonance, or bodily organization during sessions, those reports are not made methodologically irrelevant by the absence of a fully stabilized mechanism. They are not the last word, but neither are they mere noise awaiting dismissal. Practice-based research exists precisely because some phenomena are encountered first as disciplined experience, effect, or transformation before they can be adequately formalized in third-person terms \parencite{borgdorff_2012, kindon_2007}. Dejours's work on the gap between prescribed work and real work is illuminating in this respect: bodies often know more than the official language of the task can capture, and the gap between codified procedure and enacted practice is not an anomaly but a structural feature of human activity \parencite{dejours_1980}. Beacon enters the program at that level. It produces experiential and physiological material that must be investigated, not protected from investigation and not erased by the demand that only already-formalized phenomena count as legitimate data.

This does not leave all forms of knowledge on the same plane. HIT does not treat testimony, statistical modeling, physiological measurement, historical reconstruction, and computational inference as equivalent simply because they all matter. The framework is plural without being indiscriminate. It treats different forms of construction, inscription, and evidentiary articulation as differently powerful for different questions and expects them to be held in tension rather than fused into a single homogenized language. That plural structure is not a weakness in the program. It is the form demanded by an object that appears as ratio, recurrence, resonance, embodiment, interpretation, and control across multiple domains at once. The epistemic wager of HIT is therefore neither that every field already knows the same thing, nor that one privileged field can absorb all the others, but that these partial knowledges can be assembled into a more disciplined and more adequate picture than any one of them could produce alone.

\section{Science, power, and the limits of formalization}

If HIT asks different regimes of inquiry to coexist within a single framework, it must also be explicit about the institutions that regulate what counts as legitimate knowledge. Scientific institutions are not neutral windows onto truth. They are historically situated arrangements of prestige, discipline, funding, publication, and gatekeeping that shape which objects become visible, which vocabularies are tolerated, and which questions are treated as serious in the first place \parencite{foucault_1980, foucault_1990}. This matters especially for a framework such as HIT, whose central terms, including vibration, resonance, and harmony, circulate simultaneously in physics, music, medicine, spirituality, engineering, and popular discourse. The same vocabulary that can name a technically precise problem can also trigger institutional suspicion because of its proximity to diffuse or unrigorous uses. This is not a hypothetical risk. Terms such as harmony or vibration often encounter reflex dismissal in technical settings before their formal content has even been examined, precisely because the same words circulate in poorly delimited therapeutic or spiritual registers. A manuscript in this space cannot afford innocence about those conditions.

The point, however, is not to denounce science from outside. It is to describe the conditions under which formalization becomes both necessary and partial. Every act of formalization translates. It moves from bodily process to report, from report to concept, from concept to metric, from metric to data structure, and from there to model or device. Fern\'andez M\'endez and San Emeterio's account of the progressive externalization of technique provides a useful language for this chain because it shows that formalization is never simple extraction; it is reconfiguration through successive material and symbolic interfaces \parencite{fernandez_2021b}. Something is gained at each step: stability, comparability, transmissibility, repeatability. Something is also altered: texture, situational density, first-person thickness, tacit know-how. That dual movement is not a defect of method. It is the condition of method.

Seen from a Foucauldian angle, this also means that no concept is exhausted by the discipline that first stabilized it. Disciplines are historical arrangements of objects, methods, and permissions rather than eternal jurisdictions over words. A term such as entropy therefore cannot be borrowed loosely, but neither must it remain imprisoned within one canonical field of use. HIT uses such concepts only under explicit reformulation, preserving their technical histories while allowing them to illuminate broader problems of order, transformation, and constraint across domains.

HIT therefore treats formalization neither as betrayal nor as salvation. A session report translated into physiological time series becomes more measurable and more comparable, but it does not remain identical to the lived event from which it was abstracted. A harmonic descriptor translated into tensorized input for a neural model becomes experimentally tractable, but it ceases to be the phenomenon in its original setting. This is why the program needs both critical sobriety and methodological confidence at once. Without formalization there is no cumulative science; without reflexivity about formalization there is no honest account of what scientific artifacts do and do not preserve. The methodological task of HIT begins precisely there: in the disciplined handling of translation under conditions of unavoidable loss and equally unavoidable gain.

\section{An epistemology of inconsistency}

The previous section established that formalization always translates. The present section adds a more difficult point: no regime of representation, including scientific language itself, ever captures its object without remainder. HIT therefore requires what may be called an epistemology of inconsistency, by which is meant not incoherence in the pejorative sense but a disciplined recognition that knowledge is always produced across heterogeneous registers whose relations are imperfect, partial, and often unstable. Bodies, habits, speech, diagrams, models, and instruments do not simply duplicate one another. They disclose overlapping but non-identical slices of a process that exceeds any one register of capture. This is also where HIT's realism must be stated with care. Models do not gain access to the Real by fully succeeding. They become answerable to it precisely where their own explanatory aims encounter contradiction, remainder, or failure of closure. What they secure is not the Real itself delivered intact to knowledge, but the patterned trace left by resistant constraints within and across registers.

This broader view is already visible in Fern\'andez M\'endez's work on bodily automatisms, malabarism, and digitized technique, and in the work with San Emeterio on technology, language, and labor, where knowledge appears as distributed across corporal, designed, machinic, organizational, and digital forms rather than concentrated in explicit verbal articulation alone \parencite{fernandez_2021a, fernandez_2021b, fernandez_2015, sanemeterio_2017a}. Le Bret\'on and Dejours reinforce this point from different directions by insisting that the body is not merely an instrument that carries knowledge but one of the conditions through which knowledge becomes possible at all \parencite{dejours_1980, le_breton_1995, le_breton_2002}. Within HIT this matters immediately. Resonant regulation, perceptual attunement, therapeutic effect, and computational salience may be linked, but they do not become interchangeable just because they coexist. Each register preserves something the others cannot fully absorb.

Maturana's rejection of information transfer offers a biological counterpart to this same limit. In the biology of cognition, organisms do not receive information from the environment in the Shannonian sense; they are perturbed, and the significance of that perturbation depends on their own structure and history of coupling \parencite{maturana_1970, maturana_varela_1987}. This does not make communication unreal. It means that the transmissive model is insufficient once living systems are at issue. A signal does not specify its effect from the outside; the receiving system enacts that effect according to its own organization. The convergence with the present chapter is exact enough to matter. Lacan formulates the inadequacy of transparent transmission from the side of language and subjectivity; Maturana formulates it from the side of biological organization. Both interrupt the fantasy that messages pass unchanged from source to destination.

Lacan sharpens the issue further by showing that the problem is not just practical but structural. If a signifier represents a subject for another signifier, then the speaking subject is never fully present to itself in what it says; it emerges within a differential chain that exceeds intentional mastery \parencite{lacan_1966, lacan_1973, lacan_1975}. The subject is split not because language sometimes fails, but because language operates through displacement, substitution, and the circulation of desire. The formulation from Seminar XI regarding the \emph{d\'efil\'es of the signifier} is especially relevant here: language is not a neutral conduit through which pure information passes unchanged, but a constrained terrain in which what can be said is already shaped by the relation between body, drive, jouissance, and symbolic form \parencite{lacan_1973}. In this sense the Real is not an object waiting behind the symbol to be seized by more exact discourse. It is what remains outside complete symbolization and returns only as limit, insistence, shock, or bodily encounter. Formal thought can organize realities; it cannot capture the Real itself.

Derrida and Nietzsche, read at this point, do not dissolve the possibility of knowledge; they intensify the need for precision about what knowledge can be. Derrida's insistence that meaning is constituted through differential spacing rather than through simple presence makes it difficult to imagine any representational system as self-grounding or complete \parencite{derrida_1968}. Nietzsche's claim that there are no uninterpreted facts in the naive sense works similarly as a warning against mistaking inherited conceptual schemas for transparent access to the world \parencite{nietzsche_1873}. In that restricted but decisive sense, inquiry never steps outside text, mediation, or symbolization into a pure object. It constructs realities through inscriptions, translations, and interpretive cuts. Within HIT this does not cancel the distinction between better and worse inquiry, but it does prevent the framework from pretending that one language, one metric, or one institution can close the object once and for all. An epistemology of inconsistency is therefore not a retreat from rigor. It is the condition under which a plural, transdisciplinary object can be studied without forcing its heterogeneous manifestations into premature equivalence.

Once that limit has been acknowledged, the methodological question shifts. If no discourse captures its object transparently, then HIT needs instruments capable of producing structured evidence that is not reducible to natural-language self-description alone. Computational models enter at precisely that juncture: not as an escape from mediation, but as a different regime of inscription in which controlled transformations can expose relations that ordinary discourse cannot stabilize by itself.

\section{Neural networks as scientific instruments}

Once the problem of representation has been made explicit, the role of computational models can be defined more sharply. Phideus, the program's computational probe, developed more fully in Chapter 11, does not treat neural networks only as optimization targets whose value lies in surpassing a benchmark. Within HIT they also function as scientific instruments. This means that their epistemic role is closer to that of a spectrometer than to that of a black-box classifier. They receive transduced traces of physical or biological phenomena, transform those traces under controlled architectural and descriptor conditions, and produce embeddings, alignments, and retrieval geometries that can be read as evidence about structured organization within the object as it has been symbolically and technically constructed by the pipeline. The goal is not to anthropomorphize the network or to confuse representation with reality. The goal is to use controlled variation in representational space as a way of observing what kinds of structure survive, reorganize, or disappear under explicit interventions.

The atoms-to-bits-to-atoms bridge is central here. A physical event is captured through a chain of sensing and preprocessing, converted into digital form, transformed by a model, and then interpreted back into claims about the constructed object that this chain makes available. Every stage in that bridge introduces its own assumptions and possible distortions. This is why Phideus cannot justify its conclusions simply by pointing to performance numbers. The epistemic question is more demanding: what kind of claim is licensed when a descriptor systematically reorganizes latent geometry, improves retrieval under adversarial controls, or leaves causal signatures under matched conditions? The spectrometer analogy clarifies the answer. A prism does not create the wavelengths it displays, but neither does it display them without imposing an optical apparatus of its own. Likewise, a network does not invent every invariant it exposes, but it also does not expose invariants without the mediation of architecture, loss function, and transduction.

For that reason HIT distinguishes among levels of computational claim. The strongest claims are structural: recurrent invariants that survive changes of architecture, domain, and instrumentation. The next level is geometric: reproducible reorganizations of representational space that suggest systematic structure without yet fixing mechanism. The weakest level is mechanistic: precise attributions about why a given descriptor works, which usually require stronger causal contrast than a single model family can provide. Phideus operates most securely at the structural and geometric levels, and only conditionally at the mechanistic one. This tiered reading is not a concession. It is what makes a modest structural realism possible inside the program's own constructed objects. If an invariant survives changes in instrument, modality, and configuration, it earns provisional epistemic weight as a candidate invariant of the scientific object rather than as a quirk of one specific model \parencite{lakatos_1978, pearl_2009, popper_1959, kornblith_2019}.

\section{The rectification: natural harmony and perceptual harmony}

The central methodological rectification of the program follows directly from the ontological and epistemological arguments developed so far. If HIT wishes to test whether harmonic ratios have privileged explanatory status, it cannot collapse physically natural organization into perceptual or cultural transforms by default. The distinction between natural and perceptual harmony therefore becomes methodologically decisive. This rectification did not arise from abstraction alone. It became necessary because early positive effects within the program already showed that changing the representation of the same phenomenon could change what became learnable, retrievable, or experientially salient without yet deciding what kind of harmonic order had actually been privileged. A useful representation is not automatically a physically natural one. Without a prior distinction between natural and perceptual harmony, experimental success could too easily be mistaken for confirmation of the stronger theoretical claim. Natural harmony refers here to ratio structure expressed in physically linear terms, including the series relations and proportional organizations that arise from the phenomenon itself. Perceptual harmony refers to transforms that are psychophysically or musically meaningful, such as logarithmic spacing, equal temperament, or semitone-based organization. Both may matter. The methodological point is that they must not be conflated at the moment of hypothesis testing.

This distinction must first be understood at the level of HIT as a whole and only afterward at the level of each experimental implementation taken in its singularity. The program cannot allow physically natural organization and perceptually normalized organization to blur into one another under a single harmonic vocabulary, because they answer different questions. Natural harmony concerns the resonant organization that emerges from the phenomenon or medium itself. Perceptual harmony concerns the transformations through which organisms, cultures, or technical systems render that organization intelligible, usable, or musically commensurable. Both belong to the program, but they do not belong to it in the same way. The first names a claim about structured relation as it is constructed from the phenomenon under a given symbolic and instrumental regime; the second names a claim about how that organization is received, transformed, and stabilized within perceptual or cultural regimes.

Once this distinction has been stated at the general level, it can take different experimental forms without changing its methodological meaning. In Phideus, discussed more fully in Chapter 11, it appears as a controlled comparison among alternative ways of constructing the same phenomenon, so that physically grounded descriptions can be distinguished from perceptually normalized or otherwise nearby codings under comparable computational conditions. In Beacon, discussed more fully in Chapter 12, the same distinction operates first as an orientation of the setting itself: resonant devices, synthetic sound, bodily calibration, interference fields, and Lissajous-based exploration may be organized from the side of natural harmony before any explicit comparison with tempered or otherwise perceptual codings is introduced. The point is not that both branches must proceed identically, but that both must state clearly which sense of harmony is being investigated and why.

Methodological discipline begins there. Whenever the distinction can be tested through direct comparison, the comparison should be made explicit and its criteria stated in advance. Whenever the work begins instead from the design of instruments, sessions, or interpretive protocols, that orientation should also be declared without pretending that a contrast has already been performed. What must remain stable is the principle itself: natural harmony and perceptual harmony are not interchangeable simply because both can be musically, culturally, or experimentally productive. HIT only becomes legible as a coherent program if that distinction is maintained at the level of method while allowing each experimental implementation to develop its own appropriate form of evidence.

\section{Beacon's methodological arsenal}

If Phideus is the program's computational probe, Beacon, developed more fully in Chapter 12, is its experiential and physiological probe. Its methodological role is not to duplicate what computation already does, but to address classes of evidence that cannot be reduced to cross-modal retrieval or descriptor ablation. This includes qualitative reports, interviews, bodily regulation, physiological synchrony, and the structured use of sound, resonance, and calibration within live sessions. Beacon therefore belongs to the manuscript not as an illustrative side project, but as one half of the program's dual experimental architecture. It asks what forms of organization become visible when one starts from lived response and bio-signal change rather than from representational learning.

The methodological repertoire of Beacon is correspondingly plural. It includes qualitative methods such as semi-structured interviews, participant testimony, and session narratives; physiological measures such as EEG, ECG or HRV, respiration, and skin conductance; and devices such as the corazofono, which attempt to register a participant's own center of resonance in a situated way \parencite{porges_2011, porges_2017}. Cymatic calibration belongs here as well, but only if stated carefully. In Beacon, visible interference patterns are not decorative metaphors for \enquote{vibration.} They function as practical instruments for calibrating or verifying harmonic organization under controlled conditions, in continuity with the longer physical archive running from Chladni to Jenny \parencite{chladni_1787, jenny_1967}. The methodological lesson is that Beacon produces evidence in registers that are irreducible to those of Phideus while remaining comparable at the level of convergent inquiry.

This plurality does not eliminate the need for caution. Beacon remains an early-stage program. Sample sizes may be small, protocols may still be evolving, and therapeutic interpretations must remain more modest than the intensity of some observed effects might tempt one to claim. But that caution is compatible with methodological seriousness. Practice-based inquiry does not treat reports as self-validating proof; it treats them as data whose form of validity differs from randomized intervention trials and from computational retrieval benchmarks. Beacon therefore enlarges the evidence architecture of HIT by bringing first-person and bio-signal regimes into disciplined relation with the rest of the program rather than leaving them outside it.

\section{The architecture of evidence}

The methodological commitments of HIT can now be stated at their most operational level. Its architecture of evidence rests on three linked requirements: convergence, controls, and adversarial testing. No isolated result settles a claim. A computational gain, a physiological shift, a qualitative report, or a single experimental run can at most contribute one line of support, but stronger conclusions require recurrence across conditions, instruments, analytical views, and registers of evidence. This architecture is not hypothetical. Across the program, changes in representation, instrument design, and protocol already alter what becomes visible, which is precisely why cumulative standards matter. For the same reason, results matter only insofar as they survive comparison with viable alternatives rather than appearing large in isolation. Those alternatives will not always take the same form: in some settings they are matched computational baselines, non-natural transforms, or shuffled conditions; in others they are baseline sessions, protocol contrasts, rival interpretations, or counter-conditions that preserve the same surface effect while altering the setup. The common rule is that a result becomes informative when it can be situated against the strongest nearby explanation short of the target hypothesis.

The adversarial principle names the final consequence of that rule. HIT does not ask only whether a pattern appears, but whether it remains legible once the best available counter-hypothesis has been allowed to contest it. This is why pre-registered expectations matter when prediction is possible, why repeated evaluation matters when seeds or sessions can drift, and why the program maintains a tiered language of claims instead of allowing every positive result to bear the same theoretical weight. Some findings support structural plausibility, others justify geometric or mechanistic readings, and only a smaller subset begin to inform the stronger claims of the framework. The chapter ends, then, not with a result but with a discipline: theory can advance only under conditions in which unlike forms of evidence are neither collapsed into one register nor treated as incommensurable by default. Once that architecture is in place, the next question can be posed with greater precision: whether natural harmonic organization can be understood as a form of informational efficiency, and under what conditions.
